\chapter{Hierarchies}

We begin by defining a `hierarchy' indexed by a club class of ordinals.

\begin{definition}[Hierarchy]
  We say a family of sets $Z_{\alpha}$ (indexed over a class $C \ssq \OR$) is a \textbf{hierarchy} if the following are all true:
  \begin{enumerate}
    \item each $Z_{\alpha}$ is transitive
    \item the class $C$ is club in $\OR$
    \item for every $\alpha < \beta$ from $C$, $Z_{\alpha} \subset Z_{\beta}$
      \item for every $\beta$ a limit point\footnote{as $C$ is club, such $\beta$s do also lie in $C$} of $C$,
      \begin{align*}
          Z_{\beta} = \bigcup_{\alpha \in C \cap \beta} Z_{\alpha}
      \end{align*}
    We usually write
    \begin{align*}
        Z = \bigcup_{\alpha \in C} Z_{\alpha}
        = \setst{x}{\exists \alpha \in C \st x \in Z_{\alpha}}
    \end{align*}
    to denote the entire hierarchy.
  \end{enumerate}
\end{definition}

\section{The Three Main Hierarchies}

\subsection{The $V$-Hierarchy}

This is also sometimes referred to as the cumulative hierarchy.

\begin{definition}[The Sets $V_{\alpha}$]
  For each ordinal $\alpha$, define $V_{\alpha}$ by recursion as follows.
  \begin{itemize}
    \item If $\alpha = 0$, define $V_{\alpha} = 0$.
    \item If $\alpha = \beta + 1$, define $V_{\alpha} = \Powset\of{V_{\alplha}}$
    \item If $\alpha$ is a limit ordinal, define $V_{\alpha} = \bigcup_{\eta < \alpha} V_{\eta}$
  \end{itemize}
  Denote
  \begin{align*}
    V = \bigcup_{\alpha \in \OR} V_{\alpha}
  \end{align*}
\end{definition}

As we refer to this as the $V$-hierarchy or the cumulative hierarchy, you might imagine that it is, indeed, a hierarchy. We will need to show this.

\subsection{The $H$-Hierarchy}

This is also sometimes referred to as the hereditary hierarchy.

\subsection{The $L$-Hierarchy}

This is also sometimes referred to as the constructible hierarchy.
